\chapter{Autoregressive Language Models}
\label{ch10:top}

Earlier chapters spoke abstractly about distributions on strings. This chapter fixes the concrete convention used for the rest of the book when the strings come from a large language model. The main conceptual move is simple but important: the mathematical object is the induced distribution on outputs, not the neural network parameters themselves.

To make that precise, we must first settle the sample space cleanly.

\section{Tokens, Vocabulary, and End-of-Sequence}

Let $\Sigma$ denote the vocabulary of ordinary content tokens, and let
\[
  \EOS \notin \Sigma
\]
be a distinguished end-of-sequence symbol. The model emits tokens from the enlarged alphabet
\[
  \Sigma \cup \{\EOS\},
\]
but the actual finite outputs are strings in
\[
  \Sigma^*.
\]
A generated string ends when the first $\EOS$ is emitted. Thus:
\begin{itemize}
  \item histories during generation are strings in $\Sigma^*$,
  \item outputs are strings in $\Sigma^*$,
  \item and $\EOS$ is a stopping action, not a symbol that appears inside the returned output string.
\end{itemize}

This convention is the cleanest one for the later generating functions. It also matches the length PGF of Chapter~8, where length means ``number of non-$\EOS$ tokens before termination.''

Tokenization still matters. A model over characters, one over words, and one over BPE tokens define genuinely different distributions on $\Sigma^*$. All later counting, length, and entropy questions depend on that choice of atomic units.

\section{The Autoregressive Factorization}

An autoregressive language model specifies, for every history $h \in \Sigma^*$, a conditional distribution
\[
  p(\cdot \mid h) \in \Delta(\Sigma \cup \{\EOS\}).
\]
If the model outputs
\[
  w = w_1 w_2 \cdots w_n \in \Sigma^*,
\]
then its probability is
\[
  p(w) = p(\EOS \mid w_1\cdots w_n) \prod_{t=1}^n p(w_t \mid w_1\cdots w_{t-1}).
\]
The empty output has probability
\[
  p(\varepsilon) = p(\EOS \mid \varepsilon).
\]

\begin{definition}[Autoregressive language model]
An \emph{autoregressive language model} over content vocabulary $\Sigma$ is a family of conditional distributions
\[
  \mathcal{M} = \{p(\cdot \mid h)
  : h \in \Sigma^*\},
\]
where each conditional is supported on $\Sigma \cup \{\EOS\}$. The induced output distribution on $\Sigma^*$ is defined by the factorization above.
\end{definition}

The conditionals after a history already containing $\EOS$ are irrelevant, because such histories are never reached by the sampling process. One may define them arbitrarily or ignore them entirely.

This is the right abstraction for the book. The network parameters are one way to represent the family $\mathcal M$; they are not themselves the primary mathematical object.

\section{Softmax and Temperature}

In practice, modern LLMs compute the conditional distribution by applying softmax to a vector of logits:
\[
  p(\sigma \mid h) = \frac{e^{\ell_\sigma(h)}}{\sum_{\tau \in \Sigma \cup \{\EOS\}} e^{\ell_\tau(h)}}.
\]
For every finite temperature parameter $T>0$, the tempered distribution is
\[
  p_T(\sigma \mid h) = \frac{e^{\ell_\sigma(h)/T}}{\sum_{\tau} e^{\ell_\tau(h)/T}}.
\]
At $T=1$ this is the native model. For every fixed finite $T>0$, every token in $\Sigma \cup \{\EOS\}$ keeps strictly positive probability. As $T \to 0$, mass concentrates on the set of maximizers of the logits; if there are ties, the limiting distribution is supported on the tied argmax set.

Temperature will matter later when we distinguish local tokenwise rescaling from more global Gibbs-style reweighting of full strings.

\section{Transformers in One Paragraph}

A transformer is one specific neural parameterization of the history-to-logit map
\[
  h \longmapsto \ell(h).
\]
For the mathematics of this chapter, the exact architecture matters less than the existence of that map. The only reason to mention the architecture at all is that some later theorems, such as tightness results of Du et al., are proved for particular model classes rather than for arbitrary conditional families.

\section{Tightness}

Chapter~8 introduced the right probability-theoretic notion.

\begin{definition}
An autoregressive language model is \emph{tight} if the total probability of finite outputs is $1$:
\[
  \sum_{w \in \Sigma^*} p(w) = 1.
\]
Equivalently, generation halts almost surely.
\end{definition}

For a tight model, the length distribution
\[
  p_n = \Pr(|W|=n)
\]
is a genuine probability distribution on $\mathbb{Z}_{\ge 0}$.

We use the Du et al.
result as a black box.

\begin{theorem}[Du et al., black box]
Under the hypotheses of their framework, softmax transformer language models are tight.
\end{theorem}

We deliberately do not reproduce a proof sketch here. Positivity of softmax alone is not enough for a self-contained proof, and the real argument is architectural and measure-theoretic. For this chapter, the safe lesson is simply that there are important LLM classes for which properness has been proved, not merely hoped for.

\section{Two Generating Functions Attached to an LLM}

Once the output distribution is fixed, there are at least two different generating functions one might attach to it.

\paragraph{1. The length PGF.}
For a tight model, define
\[
  P(z) = \sum_{n \ge 0} p_n z^n,
  \qquad
  p_n = \Pr(|W|=n).
\]
This is the right object for questions about how long the model tends to generate before stopping.

\paragraph{2. A counting-type typical-set GF.}
If one later chooses an entropy notion and a typicality criterion as in Chapter~9, one can count how many outputs of length $n$ are ``typical'' and form a second generating function
\[
  A(z) = \sum_{n \ge 0} a_n z^n.
\]
This is a counting object, not a length-probability object.

The point of this section is not to define every such counting GF in full detail again. It is to prevent confusion: `P(z)` and `A(z)` answer different questions and generally have different radii and different singularity interpretations.

\section{Prefix Dependence and the Markov Viewpoint}

There is one useful formal viewpoint. The sampling process can be seen as a Markov chain on the countable state space of finite prefixes $\Sigma^*$. From any history there are only finitely many next-token choices, so the branching is large but still finite at each step.

However, Chapter~9's entropy-rate theory for stationary infinite processes does \emph{not} transfer automatically to this terminated setting. A tight autoregressive model produces a finite output and then stops. To speak about an entropy rate in the stationary sense, one would need an additional construction, such as an associated infinite process or a conditioning scheme. We therefore avoid claiming in this chapter that a generic tight LLM comes with a Chapter~9 entropy rate for free.

What we can say safely is simpler:
\begin{itemize}
  \item finite-context models behave more like finite-order Markov sources,
  \item full-prefix models can depend on the entire history,
  \item and the generating functions of their output distributions need not be rational unless further finite-state structure is present.
\end{itemize}

\section{A Unigram Example}

Consider the simplest autoregressive model: at every step, independently of history, emit a non-$\EOS$ token with total probability $1-q$ and emit $\EOS$ with probability $q$.

Then
\[
  \Pr(|W|=n) = q(1-q)^n
  \qquad (n \ge 0),
\]
so the length PGF is
\[
  P(z) = \sum_{n \ge 0} q(1-q)^n z^n = \frac{q}{1-(1-q)z}.
\]
This is a rational function with a simple pole at
\[
  z = \frac{1}{1-q} > 1.
\]
So even the most trivial autoregressive model already fits the rational pole picture of Part~I.

\section*{Preview}

The next two chapters ask how far that rational picture extends.

\begin{itemize}
  \item Chapter~11 studies what kinds of formal-language behavior transformer architectures can and cannot support.
  \item Chapter~12 asks under what approximation targets a finite-state weighted model is a useful surrogate for a neural language model.
\end{itemize}

Those are the chapters where ``close to rational'' must be made precise. This chapter only sets up the underlying probabilistic object and the two generating functions one might study.

\section*{Further Reading}

\begin{itemize}
  \item \textbf{A. Vaswani et al.}, ``Attention is all you need'' (2017) \cite{VaswaniEtAl2017} --- the original transformer architecture paper.
  \item \textbf{L. Du et al.}, ``A measure-theoretic characterization of tight language models'' (2023) \cite{DuEtAl2023} --- the main reference for tightness of autoregressive language models.
  \item \textbf{R. Cotterell et al.}, ``Formal aspects of language modeling'' (2023) \cite{CotterellEtAl2023} --- for a broader formal perspective on language models as probability measures on strings.
  \item \textbf{T. Icard}, ``Calibrating generative models: the probabilistic Chomsky--Schützenberger hierarchy'' (2020) \cite{Icard2020} --- for one formal hierarchy comparing rational and algebraic probabilistic models.
\end{itemize}
